\chapter{Aplicaciones y ejemplos}

%: EJEMPLOS SUCESIONES POLINOMIOS
\section{Sucesiones de polinomios conocidas}

%: POLINOMIOS BERNOULLI
\subsection{Polinomios de Bernoulli}

\begin{definicion}
    Se llaman \textbf{polinomios de Bernoulli de orden $r$} a los polinomios $B_n^{(r)}(x)$ tales que forman las sucesión de Appell para
    \begin{equation*}
        G(T) \equiv (\frac{e^T - 1}{T})^r \equiv (\frac{\Delta}{T})^r \quad (r \neq 0)
    \end{equation*}
    es decir, $B_n^{(r)}(x) \sim (T, (\frac{e^T - 1}{T})^r)$. En el caso de que $r = 1$, se llaman simplemente \textbf{polinomios de Bernoulli} y se escriben como $B_n(x)$.
\end{definicion}

\begin{definicion}
    Se llaman \textbf{números de Bernoulli de orden $r$} a los escalares $B_n^{(r)}(0)$ y se simbolizan como $B_n^{(r)}$. Si $r=1$, se llaman simplemente \textbf{números de Bernoulli} y se escriben como $B_n$.
\end{definicion}

Con esta definición, la caracterización \ref{caracterizacion diferencial} y el teorema \ref{funcion generadora appell y asociada} para sucesiones de Appell, nos proporcionan su caracterización diferencial y su función generadora, respectivamente, de forma trivial. 

\begin{proposicion}\label{recurrencia bernoulli}
    \textbf{(ecuación diferencial)}. Los polinomios de Bernoulli de orden $r$ cumplen
    \begin{equation*}
    T B_n^{(r)}(x)= B_n^{(r)}(x)' = n B_{n-1}^{(r)}(x)
    \end{equation*}
\end{proposicion}

\begin{proposicion}\label{generadora polinomios bernoulli}
    \textbf{(función generadora)}. Los polinomios de Bernoulli de orden $r$ tienen la función generadora
    \begin{equation*}
    \sum_{k=0}^\infty \frac{B_k^{(r)}(x)}{k!} T^k = (\frac{T}{e^T - 1})^r e^{x T}
    \end{equation*}
\end{proposicion}

\begin{corolario}\label{generadora numeros bernoulli}
    Los números de Bernoulli de orden $r$ tienen la función generadora
    \begin{equation*}
    \sum_{k=0}^\infty \frac{B_k^{(r)}}{k!} T^k = (\frac{T}{e^T - 1})^r
    \end{equation*}
\end{corolario}

\begin{proposicion}\label{bernoulli orden r en funcion de bernoullis}
    Los números de Bernoulli de orden $r$ verifican
    \begin{equation*}
    B_n^{(r)} = \sum_{l_1 + \cdots l_r = n}^n \binom{n}{l_1, \dots, l_r} B_{l_1} \cdots B_{l_r}
    \end{equation*}
\end{proposicion}

\begin{demo}
    Reescribiendo el lado derecho del corolario anterior, tenemos que
    \begin{equation*}
    (\frac{T}{e^T - 1})^r = \prod_{m=0}^r(\frac{T}{e^T - 1}) = \prod_{m=0}^r(\sum_{k=0}^\infty \frac{B_k}{k!} T^k) = \sum_{k=0}^\infty (\sum_{l_1 + \cdots l_r = n}^n \binom{n}{l_1, \dots, l_r} B_{l_1} \cdots B_{l_r})T^k
    \end{equation*}
    e igualando con el lado izquierdo se llega al resultado.
\end{demo}

\begin{proposicion}\label{polinomio bernoulli operador sobre x^n}
    Los polinomios de Bernoulli de orden $r$ vienen dados por
    \begin{equation*}
    B_n^{(r)}(x) = (\frac{T}{e^T - 1})^r x^n 
    \end{equation*}
\end{proposicion}

\begin{demo}
    Basta con usar \ref{relacion appell y asociada}.
\end{demo}

\begin{proposicion}\label{bernoulli r+s en terminos de bernoulli r}
    Para todo $r$ y $s$ se verifica
    \begin{equation*}
    B_n^{(r+s)}(x) = (\frac{T}{e^T - 1})^s B_n^{(r)}
    \end{equation*}
\end{proposicion}

\begin{demo}
    A partir de la proposición anterior llegamos a que
    \begin{equation*}
    B_n^{(r+s)}(x) = (\frac{T}{e^T - 1})^{(r+s)} x^n = (\frac{T}{e^T - 1})^s (\frac{T}{e^T - 1})^r x^n = (\frac{T}{e^T - 1})^s (\frac{T}{e^T - 1})^s B_n^{(r)}
    \end{equation*}
    como queríamos ver.
\end{demo}

\begin{proposicion} \label{ecuacion recurrencia bernoulli}
    \textbf{(ecuación de recurrencia)}. Para todo $r \neq 1$, los polinomios de Bernoulli vienen dados por la ecuación de recurrencia
    \begin{equation*}
        B_n^{(r)}(x +1) = B_n^{(r)}(x) + nB_{n-1}^{(r-1)}(x)
    \end{equation*}
    y para $r = 1$ vienen dados por
    \begin{equation*}
        B_n(x+1) = B_n(x) + nx^{n-1}
    \end{equation*}
\end{proposicion}

\begin{demo}
    Utilizando \ref{recurrencia bernoulli} y \ref{polinomio bernoulli operador sobre x^n} obtenemos el resultado pues
    \begin{equation*}
        B_n^{(r)}(x+1) - B_n^{(r)}(x) = \Delta B_n^{(r)}(x) = \Delta (\frac{T}{\Delta})^r x^n = (\frac{T}{\Delta})^{r-1} T x^n = nB_{n-1}^{(r-1)}(x)
    \end{equation*}
\end{demo}

\begin{corolario} \label{simetria bernoulli}
    Los polinomios de Bernoulli verifican la simetría
    \begin{equation*}
        B_n(1) = B_n
    \end{equation*}
\end{corolario}

\begin{corolario}
    Se tiene la igualdad
    \begin{equation*}
        \sum_{k=0}^n k^p =  \frac{B_{p+1}(n+1) - B_{p+1}}{p+1}
    \end{equation*}
    para todo $p \neq 0$.
\end{corolario}

\begin{demo}
    La proposoicón anterior nos lo da, ya que
    \begin{equation*}
        \sum_{k=0}^n k^p = \sum_{k=0} \frac{B_{p+1}(k+1) - B_{p+1}(k)}{p+1} = \frac{B_{p+1}(n+1) - B_{p+1}}{p+1}
    \end{equation*}
\end{demo}

\begin{proposicion}\label{recurrencia suma bernoulli}
    \textbf{(ecuación de recurrencia)}. Los polinomios de Bernoulli $B_n^{(r)}(x)$ satisfacen la ecuación de recurrencia
    \begin{equation*}
        \sum_{k=1}^n \binom{n}{k} B_{n-k}(x) = nx^{n-1} \quad (n \geq 1)
    \end{equation*}
    En particular, los números de Bernoulli cumplen que
    \begin{equation*}
        \sum_{k=1}^n \binom{n}{k} B_{n-k} = 0 \quad (n \geq 1)
    \end{equation*}
\end{proposicion}

\begin{demo}
    Por la propia definición de $G(T)$ como serie formal, tenemos que para todo entero $m \geq 1$
    \begin{equation*}
        \frac{e^T - 1}{T} = \sum_{k=1}^\infty \frac{T^k}{k!} = \sum_{k=m}^\infty \frac{T^{k-m}}{(k-m+1)!} \implies T^m \frac{e^T - 1}{T} = \sum_{k=m}^\infty \frac{T^{k}}{(k-m+1)!}
    \end{equation*}
    Y, por tanto
    \begin{equation*}
        T^m = \sum_{k=m}^\infty \frac{1}{(k-m+1)!} (\frac{T}{e^T - 1}) T^k
    \end{equation*}
    Así pues, si lo aplicamos como operadores sobre $x^n$, usando \ref{polinomio bernoulli operador sobre x^n}, vemos que
    \begin{align*}
        (n)_m x^{n-m} = \sum_{k=m}^n \frac{(n)_k}{(k-m+1)!} B_{n-k}(x)
    \end{align*}
    de donde obtenemos el resultado tomando $m=1$.
\end{demo}


La identidad de Appell \ref{identidad binomial y de Appell} nos proporciona varios resultados.
\begin{proposicion}\label{identidad appell bernoulli}
    \textbf{(identidad de Appell)}. Los polinomios de Bernoulli de orden $r$ satisfacen
    \begin{equation*} 
    B_n^{(r)}(x + y) = \sum_{k=0}^n \binom{n}{k} B_k^{(r)}(y)x^{n-k}
    \end{equation*}
\end{proposicion}

\begin{corolario}\label{polinomios bernoulli funcion numeros bernoulli}
    Los polinomios de Bernoulli de orden $r$ se pueden escribir como
    \begin{equation*}
    B_n^{(r)}(x) = \sum_{k=0}^n \binom{n}{k} B_{n-k}^{(r)}x^k
    \end{equation*}
\end{corolario}

\begin{demo}
    Tomando $y = 0$ en la proposición anterior.
\end{demo}

\begin{corolario}
    Los polinomios de Bernoulli de orden $r$ se pueden escribir como
    \begin{equation*}
     B_n^{(r+s)}(x + y) = \sum_{k=0}^n \binom{n}{k} B_k^{(r)}(y)B_{n-k}^{(s)}(y)
    \end{equation*}
\end{corolario}

\begin{demo}
    Aplicando $(\frac{T}{e^T - 1})^s$ a \ref{identidad appell bernoulli} y usando \ref{bernoulli r+s en terminos de bernoulli r} se obtiene el resultado.
\end{demo}

\begin{teorema}
    \textbf{(sumación de Euler-Maclaurin)} Dado $p(x) \in \poly$ se verifica que
    \begin{equation*}
        \sum_{i=0}^{\mu-1} p(\lambda+i) = \int_{\lambda}^{\lambda + \mu} p(u)du + \sum_{k = 1}^\infty \frac{B_k}{k!}(p^{(k-1)}(\lambda + \mu) - p^{(k-1)}(\lambda))
    \end{equation*}
    para todo $\lambda, \mu \in \field$.
\end{teorema}

\begin{demo}
    Por el teorema de expansión \ref{teorema de expansion} para los polinomios de Bernoulli $B_n(x)$, tenemos que
    \begin{equation*}
        E_y(T) = \sum_{k=0}^\infty \frac{\langle E_y(T)| B_k(x) \rangle}{k!} \frac{e^T - 1}{T} T^k = \frac{e^T - 1}{T} + \sum_{k=1}^\infty \frac{B_k(y)}{k!} (e^T - 1) T^{k-1} 
    \end{equation*}
    Si ahora lo aplicamos como operador a un polinomio genérico $p(x) \in \poly$ obtenemos
    \begin{equation*}
        p(x+y) = \int_{x}^{x+1} p(u)du + \sum_{k=1}^\infty \frac{B_k(y)}{k!} (p^{(k-1)}(x+1) - p^{(k-1)}(x))
    \end{equation*}
    Tomando $y = 0$ y sustituyendo $x$ por $\lambda, \lambda + 1, \dots, \lambda + \mu -1$, llegamos a las ecuaciones
    \begin{equation*}
        \begin{dcases*}
            p(\lambda) = \int_{\lambda}^{\lambda + 1} p(u)du + \sum_{k=1}^\infty \frac{B_k}{k!} (p^{(k-1)}(\lambda + 1) - p^{(k-1)}(\lambda)) \\
            p(\lambda + 1) = \int_{\lambda+1}^{\lambda + 2} p(u)du + \sum_{k=1}^\infty \frac{B_k}{k!} (p^{(k-1)}(\lambda + 2) - p^{(k-1)}(\lambda + 1))\\
            \vdots \\
            p(\lambda + \mu - 1) = \int_{\lambda + \mu -1}^{\lambda + \mu} p(u)du + \sum_{k=1}^\infty \frac{B_k}{k!} (p^{(k-1)}(\lambda + \mu) - p^{(k-1)}(\lambda + \mu - 1))
        \end{dcases*}
    \end{equation*}
    y sumando todas ellas obtenemos el resultado buscado.
\end{demo}

A partir de aquí obtenemos por recurrencia tanto los valores de los números de Bernoulli como el de los polinomios de Bernoulli. En particular, usando \eqref{polinomios bernoulli funcion numeros bernoulli} podemos calcular los primeros para obtener los segundos. Los primeros valores son
\begin{itemize}
    \item $B_0(x) = 1$ y $B_0 = 1$ como ya sabemos.
    \item $2B_1 + B_0 = 0 \implies B_1 = -\frac{1}{2}$ y $B_1(x) = B_1 + B_0 x \implies B_1(x) = 1 - \frac{x}{2}$
    \item $3B_2 + 3B_1 + B_0 = 0 \implies B_2 = \frac{1}{6}$ y $B_2(x) = B_2 + 2B_1 x + B_0 x^2 \implies B_2(x) = \frac{1}{6} - x + x^2$
    \item $4B_3 + 6B_2 + 4B_1 + B_0 = 0 \implies B_3 = 0$ y $B_3(x) = B_3 + 3B_2 x + 3B_1 x^2 + B_0x^3 \implies B_3(x) = \frac{1}{2}x - \frac{3}{2}x^2 + x^3$
\end{itemize}
Que $B_3 = 0$ no es ninguna casualidad. Si uno continúa calculando los números de Bernoulli, observará que todos los números en valores impares (salvo para $n=1$) se anulan. Esto es consecuencia del siguiente resultado.

\begin{teorema}
    \textbf{(fórmula del argumento complementario)}. Los polinomios de Bernoulli de orden $r$ verifican que 
    \begin{equation*}
        B_n^{(r)}(-x) = (-1)^n B_n^{(r)}(x + r)
    \end{equation*}
\end{teorema}

\begin{demo}
    Usando \ref{producto escalares} vemos que
    \begin{equation*}
        \langle (\frac{e^{-T} - 1}{-T})^r| B_n^{(r)}(-x)\rangle = \langle (\frac{e^{T} - 1}{T})^r| B_n^{(r)}(x)\rangle = n! \delta_{n,k}
    \end{equation*}
    con lo que $B_n^{(r)}(-x) \sim ((\frac{e^{-T} - 1}{-T})^r, -T)$. Ya que $(-x)^n \sim (-T)$, la proposición \ref{relacion sheffer y asociada} nos da
    \begin{align*}
        B_n^{(r)}(-x) = (\frac{-T}{e^{-T} - 1})^r (-x^n) & = e^{r T} (\frac{T}{e^T - 1})^r (-x^n) \\
        & = (-1)^n E_r(T) B_n^{(r)}(x) = (-1)^n B_n^{(r)}(x + r)
    \end{align*}
    que era lo que queríamos ver.
\end{demo}

\begin{corolario}
    Para todo entero $m \geq 1$ se verifica que $B_{2m + 1} = 0$. 
\end{corolario}

\begin{demo}
    Por el teorema anterior, tomando $r=1$ y $x=0$, y por la relación de simetria \ref{simetria bernoulli}, obtenemos que
    \begin{equation*}
        B_n = (-1)^n B_n(1) = (-1)^n B_n
    \end{equation*}
    de donde se sigue el resultado.
\end{demo}

\begin{corolario}
    Para todo entero $m \geq 1$ se verifica que $B_{2m + 1}(1) = 0$. 
\end{corolario}

\begin{proposicion}\label{caracterizacion integral bernoulli}
    \textbf{(caracterización integral)}. Los polinomios de Bernoulli de orden $r$ satisfacen la ecuación integral
    \begin{equation*}
        \int_{x+ y}^{x} B_n^{(r)}(u) du = \frac{B_{n+1}^{(r)}(x + y) - B_{n+1}^{(r)}(x)}{n+1}
    \end{equation*}
    para todo $x, y \in \field$. En particular
    \begin{equation*}
        \int_{0}^{x} B_n^{(r)}(u) du = \frac{B_{n+1}^{(r)}(x) - B_{n+1}^{(r)}}{n+1}
    \end{equation*}
\end{proposicion}

\begin{demo}
    Por la definición del operador de Bernoulli vemos que
    \begin{equation*}
        \int_{x+ y}^{x} B_n^{(r)}(u) du = (\frac{e^{yT} - 1}{T}) B_n(x) = \sum_{k=1}^{\infty} \frac{y^k}{k!} T^{k-1}B_{n}(x) 
    \end{equation*}
    Utilizando ahora la caracterizacion diferencial \ref{recurrencia bernoulli} recurrentemente en el mimebro de la derecha obtenemos que
    \begin{equation*}
        \sum_{k=1}^{\infty} \frac{y^k}{k!} T^{k-1}B_{n}(x) = \sum_{k=1}^{n+1} \frac{y^k}{k!} (n)_{k-1}B_{(n+1) - k}(x) = \frac{1}{n+1} \sum_{k=1}^{n+1} \binom{n+1}{k} y^k B_{(n+1) - k}(x)
    \end{equation*}
    de donde se sigue el resultado, por la identidad de Appell \ref{identidad appell bernoulli}.
\end{demo}

Aplicando el teorema de multiplicación para sucesiones de Appell \ref{teorema de multiplicacion} sobre valores enteros obtenemos el teorema de multiplicación para los polinomios de Bernoulli.

\begin{teorema}
    \textbf{(de multiplicación para los polinomios de Bernoulli)}. Para todo entero $m \geq 0$ se verifica
    \begin{equation}
        B_n(mx) = m^{n-1} \sum_{k=0}^{m-1} B_n(x + \frac{k}{m})
    \end{equation}
\end{teorema}

Por último, la teoría umbral aplicada sobre los números de Bernoulli puede ser aprovechada para calcular los valores de los denominados \textit{números de Stirling de segunda especie}.

\begin{definicion}
    Se definen los \textbf{k-ésimos números de Stirling de segunda especie}, ${n \brace k}$, como los coeficientes del desarrollo del operador $\Delta^k$, es decir
    \begin{equation*}
        {n \brace k} = \frac{\langle \Delta^k | x^n \rangle}{k!}
    \end{equation*}
\end{definicion}

\begin{proposicion}
    Los $k$-ésimos números de Stirling de segunda especie vienen dados por los números de Bernoulli de orden $-k$ a través de la fórmula
    \begin{equation*}
        {n \brace k} = \binom{n}{k}B_{n-k}^{(-k)}
    \end{equation*}
\end{proposicion}

\begin{demo}
    Usando \ref{polinomio bernoulli operador sobre x^n} obtenemos el resultado
    \begin{align*}
        {n \brace k} = \frac{\langle \Delta^k | x^n \rangle}{k!} = \frac{\langle \frac{\Delta^k}{T^k} T^k | x^n \rangle}{k!} = \frac{\langle (\frac{\Delta}{T})^k | T^k x^n \rangle}{k!} = \binom{n}{k} \langle 1 | (\frac{T}{\Delta})^{-k} x^{n-k} \rangle = \binom{n}{k}B_{n-k}^{(-k)}
    \end{align*}
\end{demo}

%%%%%%%%%%%%%%%%%%%%

%: POLINOMIOS LAGUERRE

\subsection{Polinomios de Laguerre}

\begin{definicion}
    Se llaman \textbf{polinomios de Laguerre de orden r} a los polinomios $L_n^{(r)}(x)$ tales que forman la sucesión de Sheffer para
    \begin{equation*}
        G(T) = \frac{1}{(1-T)^{r+1}} \quad F(T) = \frac{T}{T-1}
    \end{equation*}

    En particular, los polinomios $L_n^{(-1)}(x) \sim \frac{T}{T-1}$ se simbolizan simplemente por $L_n(x)$ y los polinomios $L_n^{0}(x)$ se denominan \textbf{polinomios simples de Laguerre}.
\end{definicion}

Es inmediato ver que
\begin{equation*}
    F'(T) = -\frac{1}{(1-T)^2} \quad \text{ y } \quad F^{-1}(T) = F(T)
\end{equation*}
Por las caracterizaciones \ref{funcion generadora} y \ref{identidad de Sheffer}, obtenemos directamente su función generadora y la identidad de Sheffer para estos polinomios.

\begin{proposicion}
    Los polinomios $L_n^{(r)}(x)$ tienen la función generadora
    \begin{equation*}
        \frac{1}{(1-T)^{r+1}}e^{\frac{xT}{T-1}} = \sum_{k=0}^\infty \frac{L_k^{(r)}(x)}{k!}T^k
    \end{equation*}
\end{proposicion}

\begin{proposicion}
    Los polinomios $L_n^{(r)}(x)$ satisfacen la igualdad
    \begin{equation*}
        L_n^{(r)}(x+y) = \sum_{k=0}^n \binom{n}{k} L_k^{(r)}(x) L_{n-k}^{(r)}(y)
    \end{equation*}
\end{proposicion}

La relación con la sucesión asociada \ref{relacion sheffer y asociada} nos da el siguiente resultado.

\begin{proposicion} \label{laguerre orden r+s funcion orden r}
    Los polinomios $L_n^{(r)}(x)$ cumplen la igualdad
    \begin{equation*}
        (1-T)^s L_n^{(r)}(x) = L_n^{(r + s)}(x)
    \end{equation*}
\end{proposicion}

\begin{demo}
    Tenemos que
    \begin{equation*}
        L_n^{(r)}(x) = (1-T)^{r+1} L_n(x)
    \end{equation*}
    Por tanto, obtenemos
    \begin{equation*}
        (1-T)^s L_n^{(r)}(x) = (1-T)^s(1-T)^{r+1} L_n(x) = (1-T)^{s+r+1} L_n(x) = L_n^{(r + s)}(x)
    \end{equation*}
\end{demo}

Usando este resultado y la caracterización \ref{caracterizacion diferencial} llegamos a ver que estos polinomios verifican la relación de recurrencia siguiente.

\begin{proposicion} \label{recurrencia diferencial laguerre}
    Los polinomios $L_n^{(r)}(x)$ verifican la ecuación de recurrencia diferencial
    \begin{equation*}
        (L_n^{(r)}(x))' = -n L_{n-1}^{(r+1)}(x)
    \end{equation*}
\end{proposicion}

\begin{demo}
    Tenemos que
    \begin{equation*}
        \frac{T}{T-1} L_n^{(r)}(x) = n L_{n-1}^{(r)}(x) \implies T L_n^{(r)}(x) = -n (1-T) L_{n-1}^{(r)}(x) = -n L_{n-1}^{(r+1)}(x)
    \end{equation*}
\end{demo}

La fórmula de recurrencia \ref{recurrencia sheffer} nos da un resultado muy útil.

\begin{proposicion}
    Los polinomios $L_n^{(r)}(x)$ cumplen la relación
    \begin{equation*}
        L_{n+1}^{(r)}(x) = (M \circ T - M + r + 1)L_n^{(r+1)}(x)
    \end{equation*}
\end{proposicion}

\begin{demo}
    La fórmula de recurrencia nos da que
    \begin{align*}
        L_{n+1}^{(r)}(x) = (M - \frac{G'(T)}{G(T)}) \circ \frac{1}{F'(T)} L_n^{(r)}(x) & = -(M - \frac{r+1}{1-T}) \circ (1-T)^2 L_n^{(r)}(x) \\
        & = (M \circ T - M + r +1)(1-T) L_n^{(r)}(x) \\
        & = (M \circ T - M + r + 1)L_n^{(r+1)}(x)
    \end{align*}
    donde la última igualdad se debe nuevamente a la proposición \ref{laguerre orden r+s funcion orden r}.
\end{demo}

A partir de este resultado, obtenemos la ecuación diferencial que estos polinomios cumplen y que suele darse como definición de los mismos.

\begin{teorema}
    \textbf{(ec. diferencial de Laguerre)}. Los polinomios $L_n^{(r)}(x)$ cumplen la ecuación diferencial
    \begin{equation*}
        xy'' + (r + 1 -x)y' + ny = 0
    \end{equation*}
\end{teorema}

\begin{demo}
    Combinando las dos proposiciones anteriores, tenemos que
    \begin{align*}
        -n L_{n-1}^{(r+1)}(x) = (L_n^{(r)}(x))' & = [(M \circ T - M + r + 1)L_n^{(r+1)}(x)]' \\
        & = [x (L_{n-1}^{(r+1)}(x))' - x L_{n-1}^{(r+1)}(x) + (r+1)L_{n-1}^{(r+1)}(x)]' \\
        & = x (L_{n-1}^{(r+1)}(x))'' + (L_{n-1}^{(r+1)}(x))' - L_{n-1}^{(r+1)}(x) 
        \\ & \quad - x (L_{n-1}^{(r+1)}(x))' + (r+1)(L_{n-1}^{(r+1)}(x))'
    \end{align*}
    de donde, reorganizando, se sigue el resultado.
\end{demo}

Una fórmula explícita de estos polinomios se puede obtener mediante las fórmulas de transferencia.

\begin{teorema}\label{formula explicita laguerre}
    Para todo $r \neq 1$, los polinomios $L_n^{(r)}(x)$ se escriben como
    \begin{equation*}
        L_n^{(r)}(x) = \sum_{k=0}^n \binom{n + r}{n - k} \frac{n!}{k!} (-x)^k
    \end{equation*}
    y en el caso en el que $r = 1$, se escriben como
    \begin{equation*}
        L_n(x) = \sum_{k=1}^{n} \binom{n-1}{k-1} \frac{n!}{k!} (-x)^{k}
    \end{equation*}
    A los números $L(n,k) = \binom{n-1}{k-1} \frac{n!}{k!}$ se les conoce como \textbf{números de Lah}.
\end{teorema}

\begin{demo}
    Veámoslo primero para $L_n(x)$. La primera de las fórmulas de transferencia \ref{transferencia asociada} nos da el resultado, pues
    \begin{align*}
        L_n(x) = -(1-T)^2(T-1)^{n+1}x^n & = (-1)^n (1-T)^{n-1} x^n \\ 
        & = (-1)^n \sum_{k=0}^{n-1} \binom{n-1}{k} (-1)^k T^k x^n = \sum_{k=0}^{n-1} \binom{n-1}{k} (-1)^{n-k} (n)_k x^{n-k} \\
        & = \sum_{k=1}^{n} \binom{n-1}{k-1} (n)_{k-1} (-x)^{n-(k-1)}
        = \sum_{k=1}^{n} \binom{n-1}{k-1} (n)_{n-k} (-x)^{k}
    \end{align*}
    donde hemos usado en la última igualdad la simetría de la suma. Para ver el caso más general, basta con usar la proposición \ref{laguerre orden r+s funcion orden r} y la caracterización \ref{sucesion asociada caracterizacion operador}, de donde vemos que
    \begin{equation*}
        L_n^{(r)}(x) = (1-T)^{r+1} L_n(x) = (-1)^n (1+T)^{r+n} x^n
    \end{equation*}
y el resultado se sigue de una forma similar a lo hecho antes.
\end{demo}

\begin{teorema}
    \textbf{(fórmula de Rodrigues)}. Los polinomios de Laguerre de orden $r$ se pueden excribir como
    \begin{equation*}
        L_n^{(r)}(x) = (M(e^x x^{-r}) \circ T^n \circ M(e^{-x})) x^{n+r}
    \end{equation*}
\end{teorema}

\begin{demo}
    Hemos visto en la demostración anterior que $L_n^{(r)}(x) = (-1)^n (1+T)^{r+n} x^n$. Además, se tiene que
    \begin{equation*}
        (M(-e^x) \circ T \circ M(e^{-x}))x^n = M(-e^x)[(Tx^{n}) \cdot e^{-x}-e^{-x} \cdot x^n] = (1 - T) x^n
    \end{equation*}
    y, por lo tanto
    \begin{equation*}
        1 - T = M(-e^x) \circ T \circ M(e^{-x}) \implies (1-T)^n = (-1)^n M(e^x) \circ T^n \circ M(e^{-x})
    \end{equation*}
    Por otro lado, se verifica además que
    \begin{align*}
        (M(x^{-r}) \circ (1-T)^n) x^{r+n} & = M(x^{-r}) \circ \sum_{k=0}^n \binom{n}{k} (-1)^k T^k x^{r+n} \\
        & = \sum_{k=0}^n \binom{n}{k} (-1)^k (r+n)_k M(x^r)x^{r+n-k}\\ 
        & = \sum_{k=0}^n \binom{n}{k} (-1)^k (r+n)_k x^{n-k}
    \end{align*}
    mientras que
    \begin{equation*}
        (1-T)^{r+n} x^n = \sum_{k=0}^\infty \binom{n+r}{k}(-1)^k T^k x^n = \sum_{k=0}^n \binom{n+r}{k}(-1)^k (n)_k x^{n-k}
    \end{equation*}
    Ahora bien, de la propia definición uno ve que $\binom{n+r}{k} (n)_k = \binom{n}{k} (n+r)_k$, con lo que se tiene que
    \begin{equation*}
        (M(x^{-r}) \circ (1-T)^n) x^{r+n} = (1-T)^{r+n} x^n 
    \end{equation*}
    Uniendo todo esto, obtenemos
    \begin{align*}
        L_n^{(r)}(x) = (-1)^n (1+T)^{r+n} x^n & = (-1)^n (M(x^{-r}) \circ (1-T)^n) x^{r+n} \\
        & = (M(x^{-r}) \circ M(e^x) \circ T^n \circ M(e^{-x})) x^{r+n} \\
        & = (M(x^{-r}e^x) \circ T^n \circ M(e^{-x})) x^{r+n} 
    \end{align*}
    que era lo que queríamos ver.
\end{demo}

%%%%%%%%%%%%%%%%%%%%%%%
%%%%%%%%%%%%%%%%%%%%%%%
%%%%%%%%%%%%%%%%%%%%%%%
%%%%%%%%%%%%%%%%%%%%%%%
%%%%%%%%%%%%%%%%%%%%%%%

%: CONEXIÓN CONSTANTES
\section{Problemas de conexión de constantes}

\begin{definicion}
    Sean $\{s_n(x)\}_{n \geq 0}$ y $\{r_n(x)\}_{n \geq 0}$ dos sucesiones de polinomios. Resolver el \textbf{problema de conexión de constantes de $s_n(x)$ con $r_n(x)$} equivale a encontrar escalares $c_{n,k} \in \field$ tales que
    \begin{equation}\label{conexion constantes}
        s_n(x) = \sum_{k=0}^n c_{n,k}r_k(x)
    \end{equation} 
    Diremos que una igualdad como esta es \textbf{la identidad $(s_n, r_n)$} y los escalares $c_{n,k}$ serán las \textbf{constantes de conexión de $s_n(x)$ con $r_n(x)$}. 
\end{definicion}

\begin{ej}
    Uno puede ver que los numeros de Stirling de segunda especie ${n \brace k}$ cumplen que
    \begin{equation*}
        x^n = \sum_{k=0}^n {n \brace k}(x)_k
    \end{equation*}
    y, por tanto, estos son las constantes de conexión que cumplen la identidad $(x^n, (x)_n)$.
\end{ej}

No es fácil encontrar una solución general al problema de conexión de constantes. Sin embargo, y para nuestra suerte, no es muy díficil ver como el cálculo umbral proporciona una solución para el caso en el que ambas sucesiones son de Sheffer.

\begin{teorema}
    Sean $s_n(x) \sim (G(T), F(T))$ y $r_n(x) \sim (H(T), L(T))$ dos sucesiones de Sheffer. Las constantes de conexión $c_{n,k}$ que resuelven el problema de conexión de constantes de $s_n(x)$ con $r_n(x)$ vienen dadas por
    \begin{equation*}
        c_{n,k} = \frac{1}{k!} \langle \frac{H(F^{-1}(T))}{G(F^{-1}(T))} L^k(F^{-1}(T))| x^n \rangle
    \end{equation*}
    Además, la sucesión de polinomios
    \begin{equation*}
        t_n(x) = \sum_{k=0}^n c_{n,k}x^n
    \end{equation*}
    es de Sheffer para $(\frac{G(F^{-1}(T))}{H(F^{-1}(T))}, F(L^{-1}(T)))$.
\end{teorema}

\begin{demo}
    Podemos reescribir \eqref{conexion constantes} como
    \begin{equation*}
        s_n(x) = (\sum_{k=0}^n c_{n,k} x^k)(\mathbf{r(x)}) = \Lambda_{H,L}\sum_{k=0}^n c_{n,k} x^k
    \end{equation*}
    y usando la estructura de grupo de los operadores de Sheffer obtenemos
    \begin{align*}
        \sum_{k=0}^n c_{n,k}x^k = \Lambda_{H, L}^{-1}s_n(x) = \Lambda_{H, L}^{-1} \circ \Lambda_{G,F}x^n & = \Lambda_{H \circ L^{-1}, L^{-1}} \circ \Lambda_{G,F}x^n \\
        & = \Lambda_{\frac{G \circ L^{-1}}{H \circ L^{-1}}, F \circ L^{-1}} x^n
    \end{align*}
    de donde vemos que $t_n(x) = \sum_{k=0}^n c_{n,k}x^k \sim (\frac{G \circ L^{-1}}{H \circ L^{-1}}, F \circ L^{-1})$. Por otro lado, usando esta igualdad junto a la adjunción de los operadores de Sheffer obtenemos
    \begin{align*}
        c_{n,k} = \frac{1}{k!} \langle T^k | \sum_{k=0}^n c_{n,k}x^k\rangle & = \frac{1}{k!} \langle T^k | \Lambda_{\frac{G \circ L^{-1}}{H \circ L^{-1}}, F \circ L^{-1}} x^n\rangle \\
        & = \frac{1}{k!} \langle \Lambda_{\frac{G \circ L^{-1}}{H \circ L^{-1}}, F \circ L^{-1}}^* T^k | x^n\rangle \\
        & = \frac{1}{k!} \langle \frac{H(F^{-1}(T))}{G(F^{-1}(T))} L^k(F^{-1}(T))| x^n \rangle
    \end{align*}
    tal y como queríamos ver. 
\end{demo}

\begin{ej}
    \begin{enumerate}
        \item Los \textbf{\textit{polinomios de Euler}} se definen como la sucesión de polinomios $E_n(x)$ que es Appell para $G(T) = \frac{e^T +1}{2}$. En particular, estos poseen la función generadora
        \begin{equation*}
            \sum_{k=0}^\infty \frac{E_k(x)}{k!}T^k = \frac{2e^{xT}}{e^T + 1}
        \end{equation*}
        Esta es bastante similar a la de los polinomios de Bernoulli que vimos en la sección anterior y es por ello natural que surja la pregunta de cómo ambos pueden relacionarse. Son varias las fórmulas que expresan uno en términos del otro; la identidad $(E_n, B_n)$ es una de ellas.

        En este caso, tenemos que $F(T) = L(T) = T$, $G(T) = \frac{e^T + 1}{2}$ y $H(T) = \frac{e^T - 1}{T}$. Esto nos da que $t_n(x)$ es de Appell para $\frac{T(e^T + 1)}{2(e^T - 1)}$. Además, la identidad de Appell nos da la igualdad
        \begin{equation}\label{appell identidad euler}
            E_n(x) = \sum_{k=0}^n \binom{n}{k} E_{n-k}(0)x^k
        \end{equation}
        Usando \ref{relacion appell y asociada}, se sigue que
        \begin{align*}
            t_n(x) =  \frac{e^T - 1}{T}\frac{2}{e^T + 1}x^n & = \frac{e^T - 1}{T} E_n(x) = \frac{e^T - 1}{T} T \frac{E_{n+1}(x)}{n+1} \\
            & = \frac{1}{n+1} (e^T - 1)E_{n+1}(x) =  \frac{1}{n+1} (e^T + 1 - 2)E_{n+1}(x) \\
            & = \frac{2}{n+1} (\frac{e^T + 1}{2}- 1)E_{n+1}(x) = \frac{2}{n+1}(x^{n+1} - E_{n+1}(x)) \\
            & = \sum_{k=0}^n \binom{n+1}{k} \frac{-2}{n+1} E_{n+1-k}(0)x^k
        \end{align*}
        donde la penúltima igualdad se debe nuevamente a \ref{relacion appell y asociada} y la última se debe a la igualdad \eqref{appell identidad euler}. Y, por tanto, obtenemos los valores de las constantes $c_{n,k}$, que nos muestran como escribir los polinomios de Euler en función de los de Bernoulli
        \begin{equation*}
            E_n(x) = \sum_{k=0}^n \binom{n+1}{k} \frac{-2}{n+1} E_{n+1-k}(0)B_k(x)
        \end{equation*}
    \item Podemos obtener una fórmula recíproca del teorema \ref{formula explicita laguerre}, es decir, podemos expresar las potencias $x^n$ en función de los polinomios de Laguerre. Para ello, basta con tener en cuenta que $x^n \sim T$ y $L_n(x) \sim \frac{T}{T-1}$, de donde vemos que los coeficientes de la identidad $(x^n, L_n(x))$ son 
        \begin{align*}
            c_{n,k} = \frac{1}{k!} \langle (\frac{T}{T-1})^k | x^n \rangle & = \frac{1}{k!} \langle T^k | (-1)^k (1-T)^{-k} x^n \rangle \\
            & = \frac{1}{k!} \langle T^k | (-1)^k (1-T)^{-k} x^n \rangle \\
            & = \frac{(-1)^k}{k!} \langle T^k | \sum_{l=0}^\infty \binom{-k}{l} (-T)^l x^n \rangle \\
            & = \frac{(-1)^k}{k!} \langle T^k | \sum_{l=0}^\infty \binom{k + l -1}{l} T^l x^n \rangle \\
            & = \frac{(-1)^k}{k!} \sum_{l=0}^\infty \binom{k + l -1}{l}\langle T^{k+l} | x^n \rangle = (-1)^k \frac{n!}{k!} \binom{n-1}{k-1}
        \end{align*}
        donde hemos usado que $\binom{-k}{l} = (-1)^k \binom{k + l -1}{l}$. Así pues, obtenemos que las potencias $x^n$ se escriben como
        \begin{equation*}
            x^n = \sum_{k=0}^{n} (-1)^k L(n,k) L_k(x)
        \end{equation*}
    \end{enumerate}
\end{ej}

%%%%%%%%%%%%%%%%%%%%%%%
%%%%%%%%%%%%%%%%%%%%%%%
%%%%%%%%%%%%%%%%%%%%%%%
%%%%%%%%%%%%%%%%%%%%%%%
%%%%%%%%%%%%%%%%%%%%%%%

%: INTERPOLACIÓN
\section{Interpolación lineal}

El cálculo umbral también permite obtener ciertas fórmulas de interpolación destinadas a construir polinomios que aproximen a una función definida en un intervalo $[a,b]$. Damos aquí tan solo una fórmula de interpolación, basada en la teoría de las sucesiones asociadas y la cual generaliza varias de las interpolaciones clásicas conocidas. No obstante, para un extenso desarrollo de esta teoría, junto con varios ejemplos, se puede consultar \cite{costabile}. 

\begin{lema}
    \textbf{(de Peano)}. Sea $\Lambda$ un operador definido sobre el anillo de funciones continuas en un intervalo $[a,b]$ tal que $\Lambda p(x) = 0$ para todo $p(x) \in \poly$ con $\gr(p(x)) \leq n$. Entonces, para toda $f \in \mathcal{C}^{n+1}[a,b]$
    \begin{equation*}
        \Lambda f(x) = \int_{a}^{b} f^{(n+1)}(t)K(t) dt
    \end{equation*}
    donde
    \begin{equation*}
        K(t) = \frac{1}{n!}\Lambda_x [(x-t)_{+}^n]
    \end{equation*}
    siendo
    \begin{equation*}
        (x-t)_{+}^n = \begin{dcases}
            (x-t)^n & x \geq t \\
            0 & x < t
        \end{dcases}
    \end{equation*}
    y donde $\Lambda_x$ indica que se piensa a $\Lambda$ como una función de $x$. 
\end{lema}

\begin{teorema}
    \textbf{(de interpolación binomial)}. Sea $G(T) \in \F^\delta$ un operador delta con sucesión asociada $p_n(x)$ y $f \in \mathcal{C}^{n+1}[a,b]G(T)$. El polinomio
    \begin{equation*}
        B_n[f](x) = \sum_{k=0}^n \frac{(G^k(T)f(x))(0)}{k!} p_k(x)
    \end{equation*}
    verifica que
    \begin{equation*}
        (G^m(T)B_n[f](x))(0) = (G^m(T)f(x))(0) \text{ y } f(x) = B_n[f](x) + R_n[f](x)
    \end{equation*}
    donde $0 \leq m \leq n$ y donde el resto $R_n[f](x)$ está dado por
    \begin{align*}
        R_n[f](x) & = \int_{a}^{b} f^{(n+1)}(t)K_n(x,t) dt\\
        K_n(x,t) & = \frac{1}{n!}[(x-t)_{+}^n - \sum_{k=0}^n \frac{[G^k(T)(x-t)_{+}^n](0)}{k!} p_k(x)]
    \end{align*}
\end{teorema}

\begin{demo}
    Tenemos que
    \begin{equation*}
        G^m(T)B_n[f](x) = \sum_{k=0}^n \frac{(G^k(T)f(x))(0)}{k!} G^m(T)p_k(x) = \sum_{k=0}^m \frac{(G^k(T)f(x))(0)}{k!} (k)_m p_{k-m}(x)
    \end{equation*}
    de donde obtenemos
    \begin{equation*}
        (G^m(T)B_n[f](x))(0) = \sum_{k=0}^m \frac{(G^k(T)f(x))(0)}{k!} (k)_m p_{k-m}(0) = (G^m(T)f(x))(0)
    \end{equation*}
    Por otro lado, $R_n g(x) \equiv R_n[g](x) = g(x) - B_n[g](x)$ cumple las condiciones del lema de Peano y de ahí obtenemos la segunda igualdad. 
\end{demo}

     

\begin{ej}
    \begin{enumerate}
        \item Cuando $G(T) = T$ y $p_n(x) = x^n$ obtenemos
        \begin{equation*}
            B_n[f](x) = \sum_{k=0}^n \frac{(T^k f(x))(0)}{k!}x^n = \sum_{k=0}^n \frac{f^{(k)}(0)}{k!}x^n
        \end{equation*} 
        y además se tendrá que
        \begin{equation*}
            B_n[f]^{(k)}(0) = f^{(k)}(0)
        \end{equation*}
        que es el clásico \textbf{teorema de Taylor}.
        \item En el caso de los factoriales descendentes, $p_n(x) = (x)_n$, estos están asociados a $G(T) = e^T -1 \equiv \Delta$ (puesto que $(e^T -1)(x)_n = n(x)_{n-1}$). Así pues, el polinomio de interpolación asociado es
        \begin{equation*}
            B_n[f](x) = \sum_{k=0}^n \frac{(\Delta^k f(x))(0)}{k!} (x)_k
        \end{equation*}
        que verifica además que $(\Delta^k B_n[f](x))(0) = (\Delta^k f(x))(0)$, de donde se deduce que
        \begin{equation*}
            B_n[f](k) = f(k) \quad k=0,1,2, \dots ,n
        \end{equation*}
        Esto no es más que la \textbf{interpolación de Newton-Gregory.}
        \item Los polinomios de Abel $A_n(x) = x(x-n)^{n-1}$ están asociados al operador de Abel $G(T) = Te^T$. El polinomio de interpolación asociado es
        \begin{equation*}
            B_n[f](x) = \sum_{k=0}^n \frac{f^{(k)}(k)}{k!} A_k(x) = \sum_{k=0}^n \frac{f^{(k)}(k)}{k!}x(x-k)^{k-1}
        \end{equation*}
        y este verificará además que
        \begin{equation*}
            B_n[f]^{(k)}(k) = f^{(k)}(k) \quad k = 0, 1, 2, \dots, n
        \end{equation*} 
        Esto es un tipo de \textbf{interpolación de Birkhoff}.
        \item Los \textit{\textbf{polinomios de Gould}} son los polinomios $p_n(x) = (\frac{x}{x - \lambda n})(\frac{x- \lambda n}{\mu})_n$ que están asociados al operador $G(T) = e^{\lambda T} (e^{\mu T} - 1) = E_{\lambda} \Delta_{\mu}$. En este caso tenemos que
        \begin{align*}
             G^k(T) f(x) & = (E_{\lambda} \Delta_{\mu})^k f(x) = \Delta_{\mu}^k E_{k\lambda} f(x) = \Delta_{\mu}^k f(x + k \lambda)\\
            & \implies (G^k(T) f(x))(0) = \Delta_{\mu}^k f(k \lambda)
        \end{align*}
        y, por tanto, el polinomio de interpolación asociado es
        \begin{equation*}
            B_n[f](x) = \sum_{k=0}^n \frac{\Delta_{\mu}^k f(k \lambda)}{k!} (\frac{x}{x - \lambda n})(\frac{x- \lambda n}{\mu})_n 
        \end{equation*}
        el cual verificará además que
        \begin{equation*}
            \Delta_{\mu}^k f(k \lambda) = \Delta_{\mu}^k B_n[f](k \lambda)
        \end{equation*}
        En particular, si tomamos $\lambda = \mu/2$ obtenemos que los polinomios de Gould para estos valores son los \textit{\textbf{factoriales centrales}} $x^{[n]}$
        \begin{equation*}
            x^{[n]} = (\frac{x}{x - \frac{\mu}{2} n})(\frac{x- \frac{\mu}{2} n}{\mu})_n = x (x + \frac{n}{2} - 1)_{n-1}
        \end{equation*}
        los cuales nos darán el polinomio de interpolación 
        \begin{equation*}
            B_n[f](x) = \sum_{k=0}^n \frac{\Delta_{\mu}^k f(-\frac{k \mu}{2})}{k!} x^{[k]}
        \end{equation*}
    \end{enumerate} 

\end{ej}

%%%%%%%%%%%%%%%%%%%%%%%
%%%%%%%%%%%%%%%%%%%%%%%
%%%%%%%%%%%%%%%%%%%%%%%
%%%%%%%%%%%%%%%%%%%%%%%
%%%%%%%%%%%%%%%%%%%%%%%